% GENERATED FILE. Edit canonical lessons, then run npm run generate:books.
% canonical-source-hash: fnv1a64:160b84c55933e823
% canonical-lessons: FR-C405-depanner, FR-C405-depanneuse, FR-C405-depanneur, FR-C405-remorquer, FR-C405-garagiste

\chapter{La dépanneuse arrive}
\label{ch:fr-la-depanneuse-arrive}
\hlchaptermodality{405}

\begin{chapteropening}
\textbf{By the end of this chapter:} \emph{I can get help after a breakdown: the tow truck, the mechanic who comes, fixing the car on the spot, towing and the garage owner.}

Complete the last lesson of chapter 405: I can get help after a breakdown: the tow truck, the mechanic who comes, fixing the car on the spot, towing and the garage owner.
\end{chapteropening}

\section[dépanner]{dépanner — to get going again (after a breakdown); to help out}

\label{lesson:FR-C405-depanner}



\begin{quote}
Before the new one: say the French for the jack, then the French for to pump up.
\end{quote}

\subsection*{You'll want to know: dépanner}
\textbf{dépanner} — \textquotedblleft{}to get going again (after a breakdown); to help out\textquotedblright{}.

\begin{quote}
Un voisin nous a dépannés avec ses câbles.
\end{quote}

\begin{quote}
\emph{A neighbour got us going again with his cables.}
\end{quote}

Also: \emph{tu peux me dépanner de dix euros ?} — can you lend me ten euros?

\begin{cousinweb}
\emph{dépanner} is built from \emph{dé-} 'un-' and \emph{panne} 'breakdown' (EtymDB). \emph{Panne} is usually traced to a sailing phrase, \emph{mettre en panne}, for a ship held still with its sails turned against the wind.
\end{cousinweb}

\begin{grammarlens}[title={What you've built}]
One more word for telling, explaining and describing at B1.
\end{grammarlens}

\subsection*{Your turn}
\begin{itemize}
\raggedright
  \item \emph{Say it:} \emph{dépanner}
  \item \emph{Say it:} \emph{dépanner}, once more
  \item \emph{Say it:} a sentence of your own with \emph{dépanner}
  \item \emph{Recall:} say the French for the jack, then the French for to pump up, then say \emph{dépanner} again
\end{itemize}

\begin{tcolorbox}[breakable,colback=teal!4,colframe=teal!35!black,title={Before you move on}]
What does dépanner mean? (\textquotedblleft{}To get going again (after a breakdown); to help out\textquotedblright{}.) Say it once more.
\end{tcolorbox}
\section[la dépanneuse]{la dépanneuse — the tow truck, the breakdown lorry}

\label{lesson:FR-C405-depanneuse}



\begin{quote}
Before the new one: say the French for to pump up, then the French for to get going again.
\end{quote}

\subsection*{You'll want to know: la dépanneuse}
\textbf{la dépanneuse} — \textquotedblleft{}the tow truck, the breakdown lorry\textquotedblright{}.

\begin{quote}
La dépanneuse est arrivée après une heure.
\end{quote}

\begin{quote}
\emph{The tow truck arrived after an hour.}
\end{quote}

Built on \emph{dépanner} with the ending \emph{-euse}.

\begin{cousinweb}
\emph{dépanneuse} is built on \emph{dépanner} 'to get going again' with the ending \emph{-euse}. \emph{Dépanner} is built from \emph{dé-} and \emph{panne} 'breakdown' (EtymDB).
\end{cousinweb}

\begin{grammarlens}[title={What you've built}]
One more word for telling, explaining and describing at B1.
\end{grammarlens}

\subsection*{Your turn}
\begin{itemize}
\raggedright
  \item \emph{Say it:} \emph{la dépanneuse}
  \item \emph{Say it:} \emph{la dépanneuse}, once more
  \item \emph{Say it:} a sentence of your own with \emph{la dépanneuse}
  \item \emph{Recall:} say the French for to pump up, then the French for to get going again, then say \emph{la dépanneuse} again
\end{itemize}

\begin{tcolorbox}[breakable,colback=teal!4,colframe=teal!35!black,title={Before you move on}]
What does la dépanneuse mean? (\textquotedblleft{}The tow truck, the breakdown lorry\textquotedblright{}.) Say it once more.
\end{tcolorbox}
\section[le dépanneur]{le dépanneur — the breakdown mechanic}

\label{lesson:FR-C405-depanneur}



\begin{quote}
Before the new one: say the French for to get going again, then the French for the tow truck.
\end{quote}

\subsection*{You'll want to know: le dépanneur}
\textbf{le dépanneur} — \textquotedblleft{}the breakdown mechanic\textquotedblright{}.

\begin{quote}
Le dépanneur a changé la batterie au bord de la route.
\end{quote}

\begin{quote}
\emph{The breakdown mechanic changed the battery by the side of the road.}
\end{quote}

In Québec, \emph{un dépanneur} is also a corner shop that stays open late.

\begin{cousinweb}
\emph{dépanneur} is built on \emph{dépanner} with the ending \emph{-eur}: the person who gets you out of a \emph{panne}. \emph{Dépanner} is built from \emph{dé-} and \emph{panne} (EtymDB).
\end{cousinweb}

\begin{grammarlens}[title={What you've built}]
One more word for telling, explaining and describing at B1.
\end{grammarlens}

\subsection*{Your turn}
\begin{itemize}
\raggedright
  \item \emph{Say it:} \emph{le dépanneur}
  \item \emph{Say it:} \emph{le dépanneur}, once more
  \item \emph{Say it:} a sentence of your own with \emph{le dépanneur}
  \item \emph{Recall:} say the French for to get going again, then the French for the tow truck, then say \emph{le dépanneur} again
\end{itemize}

\begin{tcolorbox}[breakable,colback=teal!4,colframe=teal!35!black,title={Before you move on}]
What does le dépanneur mean? (\textquotedblleft{}The breakdown mechanic\textquotedblright{}.) Say it once more.
\end{tcolorbox}
\section[remorquer]{remorquer — to tow}

\label{lesson:FR-C405-remorquer}



\begin{quote}
Before the new one: say the French for the tow truck, then the French for the breakdown mechanic.
\end{quote}

\subsection*{You'll want to know: remorquer}
\textbf{remorquer} — \textquotedblleft{}to tow\textquotedblright{}.

\begin{quote}
Ils ont remorqué la voiture jusqu'au garage.
\end{quote}

\begin{quote}
\emph{They towed the car to the garage.}
\end{quote}

\emph{Une remorque} — a trailer.

\begin{cousinweb}
\emph{remorquer} was borrowed from Italian \emph{rimorchiare} 'to tow' (EtymDB), which is usually traced through Late Latin to a Greek word for towing with a rope.
\end{cousinweb}

\begin{grammarlens}[title={What you've built}]
One more word for telling, explaining and describing at B1.
\end{grammarlens}

\subsection*{Your turn}
\begin{itemize}
\raggedright
  \item \emph{Say it:} \emph{remorquer}
  \item \emph{Say it:} \emph{remorquer}, once more
  \item \emph{Say it:} a sentence of your own with \emph{remorquer}
  \item \emph{Recall:} say the French for the tow truck, then the French for the breakdown mechanic, then say \emph{remorquer} again
\end{itemize}

\begin{tcolorbox}[breakable,colback=teal!4,colframe=teal!35!black,title={Before you move on}]
What does remorquer mean? (\textquotedblleft{}To tow\textquotedblright{}.) Say it once more.
\end{tcolorbox}
\section[le garagiste]{le garagiste — the garage owner, the car mechanic}

\label{lesson:FR-C405-garagiste}



\begin{quote}
Before the new one: say the French for the breakdown mechanic, then the French for to tow.
\end{quote}

\subsection*{You'll want to know: le garagiste}
\textbf{le garagiste} — \textquotedblleft{}the garage owner, the car mechanic\textquotedblright{}.

\begin{quote}
Le garagiste a trouvé le problème tout de suite.
\end{quote}

\begin{quote}
\emph{The mechanic found the problem straight away.}
\end{quote}

\emph{Le garagiste} runs the garage; \emph{le mécanicien} does the repairs.

\begin{cousinweb}
\emph{garagiste} is built on \emph{garage} with \emph{-iste}. \emph{Garage} comes from \emph{garer} 'to put in a safe place', which is usually traced to a Germanic verb meaning 'to watch over'.
\end{cousinweb}

\begin{grammarlens}[title={What you've built}]
One more word for telling, explaining and describing at B1.
\end{grammarlens}

\subsection*{Your turn}
\begin{itemize}
\raggedright
  \item \emph{Say it:} \emph{le garagiste}
  \item \emph{Say it:} \emph{le garagiste}, once more
  \item \emph{Say it:} a sentence of your own with \emph{le garagiste}
  \item \emph{Recall:} say the French for the breakdown mechanic, then the French for to tow, then say \emph{le garagiste} again
\end{itemize}

\begin{tcolorbox}[breakable,colback=teal!4,colframe=teal!35!black,title={Before you move on}]
What does le garagiste mean? (\textquotedblleft{}The garage owner, the car mechanic\textquotedblright{}.) Say it once more.
\end{tcolorbox}
